\subsection{A mixed two-cell example for full torsion saturation}

\begin{proposition}[A-L09: Proved; full torsion-$H^1$ saturation can change horizontal homology]
\label{AL09}
Fix integers $m>1$ and $w\geq3$, and put
$S=\Gamma_{\mathbb Z}(T)$.
There is a free integral simplicial strict PD differential
$S$-algebra $X$, obtained from the standard theory $G$ by
two attachments of AT14, with the following properties.
Its simplex-zero algebra is $S$, its support satisfies
$q\leq w'$ in every weight $w'$, its horizontal rows
retract onto those of $G$, and the coefficient map
identifies $H^0(X_n)=S$ and is a rational quasi-isomorphism
in every simplex.  Both specified target maps are the
projection $X\to G$.

Apply the full single-weight torsion-$H^1$ saturation
of A-S09-14 in weight $w$.  The resulting form-zero
row has a direct summand $\mathbb Z/m$ in normalized
homological degree one.  Thus that saturation does
not preserve exactness of this row, although both
target maps extend.  The strict PD subalgebra of
$X\otimes\mathbb Q$ generated by the saturated
weight-$w$ zero-forms has the same failure in weight $w$.

If the linear saturation is performed in every
weight while leaving all positive-form modules
unchanged, its form-zero row has a direct summand
$\mathbb Z/m$ in homological degree one in every
weight at least $w$.  This last assertion is about
the linear complex, not its later PD closure.
\end{proposition}

\begin{proof}
Write $P=D(2)$ and $Q=D(3)$ for the free contractible
simplicial abelian groups used in AT14.
Their normalized complexes have bases
\[
 P:\quad \mathbb Z e_2\xrightarrow{1}\mathbb Z e_1,
 \qquad
 Q:\quad \mathbb Z f_3\xrightarrow{1}\mathbb Z f_2.
\]
Let $\beta:P\to Q$ be the simplicial linear map
whose normalized map sends $e_2$ to $f_2$ and
$e_1$ to zero.  These formulas commute with the
chain differentials, and Dold--Kan specifies the
actual simplicial map.

Over $G$, adjoin copies $v(P),a(P),b(Q),c(Q)$ in
form degrees $0,1,1,2$, respectively, all of weight
$w$.  The underlying algebra is the free strict
PD extension, with even divided-power generators
$v,c$ and odd exterior generators $a,b$.
Define its differential, for $s\in P$ and $t\in Q$, by
\[
 \begin{aligned}
 dv(s)&=m a(s)-b(\beta s),& da(s)&=c(\beta s),\\
 db(t)&=m c(t),&dc(t)&=0.
 \end{aligned}
 \tag{AL.45}
\]
It commutes with simplicial maps.  The only nontrivial
square check is
$d^2v(s)=m c(\beta s)-m c(\beta s)=0$.
The free PD/exterior differentiation rule therefore
defines the strict PD DGA $X$.
Projection killing all four generator copies is an
actual coefficient-preserving strict PD cochain map
$X\to G$, and supplies both target maps.

More specifically, first apply AT14 with $q=3$,
parameter $Q$, zero cycle and zero old primitive;
its two generators are $c,b$ with $db=mc$.
Next apply AT14 with $q=2$, parameter $P$, cycle
$c\beta$, and old primitive $b\beta$ satisfying
$d(b\beta)=m(c\beta)$; its generators are $a,v$.
Thus \textup{(AL.45)} is precisely a pair of the
specified attachments, not an arbitrary replacement
of their module complexes.
All target generator images can be chosen zero.

The underlying parameter $P\oplus Q$ is contractible.
Applying its additive simplicial contraction to the
free PD algebra, with a degeneracy on the $G$ factor,
gives the row retractions exactly as in AT14.
The terms are free abelian groups.  Since
$P_0=Q_0=0$, we have $X_0=G_0=S$.
The assigned generator degrees are at most their
weights, so the support statement follows by adding
degrees and weights in every monomial.

Over $\mathbb Q$, replace $a(s)$ by
$a'(s)=a(s)-b(\beta s)/m$.  This is an invertible
linear change of the free odd generators, with
$G,v,b,c$ fixed.  Its differential is
\[
 dv=m a',\quad da'=0,\qquad db=m c,\quad dc=0.
\]
After rescaling by the rational unit $m$, these
are independent even-to-odd and odd-to-even
identity disks.  Their positive arities contract
over $\mathbb Q$, proving the rational coefficient
quasi-isomorphism.  A closed zero-form of $X_n$
is consequently an actual rational coefficient:
there are no negative-form boundaries.
Its image under $X\to G$ is an integral closed
zero-form, hence an element of $S$.
The coefficient-preserving map and injectivity of
$X\to X\otimes\mathbb Q$ show that the original
coefficient was integral.  Thus $H^0(X_n)=S$.
All positive vertical cohomology is torsion.

Give every new generator arity one and $G$ arity
zero.  All maps and the differential preserve
this extra arity.  In arity one the module of
new variables is the cochain complex
\[
 V^0=P\xrightarrow{\ (m,-\beta)\ }
 V^1=P\oplus Q\xrightarrow{\ (\beta,m)\ }V^2=Q.
 \tag{AL.46}
\]
Normalization in the simplicial direction is exact,
so it commutes with its vertical cohomology.
The complete nonzero normalized vertical complexes
are as follows:
\[
\begin{array}{c|c}
\text{homological degree}&\text{vertical complex}\\\hline
1&\mathbb Z\xrightarrow{m}\mathbb Z\longrightarrow0\\
2&\mathbb Z\xrightarrow{(m,-1)}\mathbb Z^2
                         \xrightarrow{(1,m)}\mathbb Z\\
3&0\longrightarrow\mathbb Z\xrightarrow{m}\mathbb Z.
\end{array}
 \tag{AL.47}
\]
The middle complex is exact: its first vector is
primitive, the last map is surjective, and its
kernel is generated by $(-m,1)$.
The first and last complexes give
\[
 H^1(V)=K(\mathbb Z/m,1),\qquad
 H^2(V)=K(\mathbb Z/m,3),\qquad H^0(V)=0.
 \tag{AL.48}
\]
These are isomorphisms of simplicial modules, specified
by the displayed normalized complexes.

In weight $w$, no monomial with two new generators
can occur.  The arity-zero part is $G^w$ and the
arity-one part has its coefficient in weight zero,
where $G^0=\mathbb Z$ in form degree zero.
The arity-one form-zero saturation can therefore be
read directly from \textup{(AL.47)}.
In normalized degree one its lattice is
$\mathbb Z(v(e_1)/m)$, because $dv(e_1)=m a(e_1)$.
In normalized degree two its lattice remains
$\mathbb Z v(e_2)$, because
$dv(e_2)=m a(e_2)-b(f_2)$ has a coefficient $-1$.
The normalized horizontal differential is now
\[
 \mathbb Z v(e_2)\xrightarrow{\ m\ }
                \mathbb Z(v(e_1)/m).
 \tag{AL.49}
\]
It has homology $\mathbb Z/m$ in degree one.
Arity zero is unchanged: $G$ has no positive
vertical cohomology, so its saturation adds no
zero-form.  This proves the direct summand in the
specified weight.  Products with a new root and
positive-weight elements, and higher divided powers
of a new root, have weight strictly greater than
$w$.  The subsequent PD closure consequently has
exactly this same weight-$w$ row.  Both target maps
extend by zero on its new arity-one roots, or by
the actual rational coefficient-preserving projection.

For completeness, the all-weight linear assertion
has a separate exact proof.  In any weight $W$, the
map sending a saturated zero-form to the vertical
class of its differential gives a short exact
sequence of simplicial abelian groups
\[
 0\longrightarrow X^{W,0}\longrightarrow
 \widetilde X^{W,0}\longrightarrow H^1(X)^W
 \longrightarrow0.
 \tag{AL.50}
\]
Here all $H^1$ is torsion by the rational calculation.
Surjectivity and the kernel assertion are exactly
the denominator and based-constant calculation in
A-S09-14.  The original row is contractible for
$W>0$, so its long exact homology sequence identifies
the horizontal homology of the middle and last terms.
The arity-one complex of $X$ is $G\otimes_S(S\otimes V)$,
with the weight shift $w$.  The coefficient inclusion
\[
 S\otimes V\longrightarrow G\otimes_S(S\otimes V)
\]
is a natural simplexwise vertical quasi-isomorphism:
the $S$-linear vertical contraction of each $G_n$
proves this, and the displayed inclusion itself
commutes with simplicial maps.
Thus \textup{(AL.48)} gives the arity-one summand
$S/mS\langle w\rangle\otimes K(\mathbb Z,1)$
of $H^1(X)$.
In every weight $W\geq w$ its horizontal degree-one
homology is $\mathbb Z/m$.
Arity supplies a genuine direct summand throughout,
which proves the final assertion via \textup{(AL.50)}.
Later multiplication by newly adjoined zero-forms can
change positive-form lattices in higher weights;
no assertion about those later PD-closed rows is
deduced from \textup{(AL.50)}.
\end{proof}

The result does not obstruct an alternating construction
which subsequently repairs the horizontal rows.
It identifies a property that the full saturation step
itself does not preserve, even for a mixed pair of the
actual AT14 attachments.
